\section*{Chapitre \ref{ch:ml:why-financial-machine-learning-is-different} --- Pourquoi l'apprentissage automatique financier est différent}

\begin{solution}{exo:ml:why-financial-machine-learning-is-different:1}
$R^2 = \rho^2 = 0.25~\%$. Pour $c = kc^\star$, $R^2 = \rho^2(2k - k^2)$~; pour $k = 3$, $-3\rho^2 = -0.75~\%$.
\end{solution}

\begin{solution}{exo:ml:why-financial-machine-learning-is-different:2}
$500/(1 + 499\times0.27) = 3.7$~; avec 0.05, $500/(1 + 24.95) = 19.3$.
\end{solution}

\begin{solution}{exo:ml:why-financial-machine-learning-is-different:3}
$(2\times0.8/0.5)^2 = 10.2$ mois~; pour trois erreurs types, $(3\times0.8/0.5)^2 = 23.0$.
\end{solution}

\begin{solution}{exo:ml:why-financial-machine-learning-is-different:4}
Le plafond est le R² du processus vrai~; un modèle ajusté au même échantillon peut aussi ajuster le bruit de cet échantillon, ce que
le processus vrai ne fait pas, de sorte que son R² dans l'échantillon peut dépasser le plafond. L'excédent (0.73 points) est du bruit ajusté~: il est
le terme de variance de la \cref{prop:ml:why:bv} vu de l'intérieur de l'échantillon, et il ne se reproduira pas. Hors échantillon, les arbres
ne peuvent pas battre 0.60~\% en espérance et obtiennent 0.57~\%.
\end{solution}

\begin{solution}{exo:ml:why-financial-machine-learning-is-different:5}
$0.045/\sqrt{240} = 0.29~\%$~: la moyenne de la fenêtre, $-0.56~\%$, est à 1.9 erreurs types du vrai zéro. Une erreur constante $e$ ajoute
$e^2$ à l'erreur quadratique moyenne~: $0.0056^2/0.0093 = 0.34$ points de R². Les rendements de la décennie de test ont une moyenne de
$+0.29~\%$, donc le terme croisé $2\times0.0056\times0.0029/0.0093 = 0.35$ points double le coût~: soit environ les 0.67 points qui séparent
la régression ridge (0.38~\%) de $-0.29~\%$.
\end{solution}

\begin{solution}{exo:ml:why-financial-machine-learning-is-different:6}
Un modèle de rendements bruts apprend une ordonnée à l'origine égale au rendement moyen de la fenêtre d'entraînement, soit la trajectoire du marché,
estimée à partir de 240 tirages~; cette seule erreur peut rendre le R² négatif (exercice~5). Entraîner le modèle sur les rendements diminués de la moyenne en coupe
transversale de chaque mois, ou évaluer la coupe séparément~: les mêmes caractéristiques donnent alors 0.38--0.57~\%.
\end{solution}

\begin{solution}{exo:ml:why-financial-machine-learning-is-different:7}
\texttt{strong\_signal(snr=0.0101/0.9899)}~: Bayes 1.01~\%~; régression ridge 0.67~\%, arbres boostés 0.42~\% (5.1~\% dans l'échantillon), réseau
0.12~\%. La fonction de Friedman est non linéaire, donc la régression ridge est biaisée, mais avec 20\,000 exemples et 1~\% de signal la variance
des modèles souples coûte plus que le biais de la régression ridge~: à faible rapport signal sur bruit, le modèle rigide l'emporte tant que l'échantillon est
suffisamment grand (les arbres du panel d'actions ont eu besoin de huit ans de données sur 500 actions).
\end{solution}

\begin{solution}{exo:ml:why-financial-machine-learning-is-different:8}
Développer $\E(y - cf)^2 = \sigma_y^2 - 2c\,\Cov(y, f) + c^2\sigma_f^2$ avec $\Cov = \rho\sigma_y\sigma_f$~; le trinôme en
$c$ atteint son maximum en $c^\star = \rho\sigma_y/\sigma_f$, avec la valeur $\rho^2\sigma_y^2$, et $R^2(2c^\star) = 0$. Une prévision qui
est corrélée à 0.92 avec le processus vrai l'est avec les rendements à environ $0.92\times\sqrt{0.006} = 0.071$~; si son échelle est
supérieure au double de $c^\star$ (un modèle ajusté à des rendements bruts portant une ordonnée à l'origine parasite, ou dont les prédictions s'étalent comme
si le signal était fort), $R^2 < 0$ bien que le classement des actions soit bon. Le coefficient d'information, qui ignore
l'échelle, est la meilleure mesure du classement~; le R² juge aussi la taille.
\end{solution}

\begin{solution}{pb:ml:why-financial-machine-learning-is-different:1}
\begin{enumerate}
\item 0.72~\% pendant les années d'entraînement, 0.60~\% pendant les années de test~: le plafond est une statistique d'échantillon du processus vrai comparée
aux rendements réalisés, et le bruit (de marché et spécifique) varie d'une décennie à l'autre.
\item $\sqrt{0.0060} = 0.077$.
\item La moyenne historique du rendement d'une action est si bruitée qu'il est facile de la battre~; zéro est l'indice de référence plus exigeant (Gu,
Kelly et Xiu).
\item 19~\% des mois (23 sur 120)~; les arbres perdent 18~\%.
\item Régression ridge 0.39~\% et 0.38~\%~; arbres 1.45~\% et 0.57~\%~; réseau 0.64~\% et 0.28~\%.
\item 64~\%, 95~\% et 47~\%.
\item Sur la tâche forte, le bruit est faible et le biais domine, donc les modèles souples gagnent et obtiennent le même score dans
l'échantillon et hors échantillon~; sur les actions, le bruit domine et la variance tranche.
\item Le terme de variance~: les arbres ont ajusté le bruit d'entraînement.
\item 3.7~: les facteurs de marché et sectoriels font que les chocs des actions sont corrélés à 0.27 en moyenne.
\item $-0.29~\%$, $-0.09~\%$, $-0.30~\%$~: leurs ordonnées à l'origine égalent le rendement moyen de la fenêtre d'entraînement, $-0.56~\%$.
\item Entre 48 mois (arbres 0.29~\%, régression ridge 0.33~\%) et 96 (0.47~\% contre 0.37~\%).
\item 4.3 mois pour s'écarter de zéro~; 21 mois pour s'écarter de la régression ridge.
\item Son R² tombe à 0.01~\% en moyenne pendant les années 6 à 10, alors que le plafond est en moyenne de 0.83~\%~; les variables, les rangs,
ont la même allure avant et après.
\item \emph{Résultat nommé}~: plafond 0.60~\%, les arbres atteignent 95~\% de celui-ci (0.57~\%), 4.3 mois pour les distinguer de zéro et 21
mois pour les distinguer de la régression ridge.
\item Avec dix ans, les arbres (avec un historique en direct assez long pour les confirmer)~; avec deux ans, la régression ridge (0.29~\% contre
0.18~\%).
\item Le R² réalisé ou l'IC, avec un test séquentiel, car ce sont les erreurs, et non les variables, qui ont révélé le changement.
\item Comme \texttt{decay\_from}~: l'effet d'une caractéristique qui s'amenuise une fois qu'elle est largement négociée.
\item Le processus vrai n'est pas observé~; il faut l'encadrer par le meilleur R² ou IC publié dans des travaux comparables, et par la
dispersion des modèles qui convergent vers un même niveau.
\item C'est une borne inférieure du plafond réel~; le fait que des modèles très différents se regroupent entre 0.33~\% et 0.40~\% suggère que le
plafond n'est pas très loin au-dessus.
\item Parce que le signal ne représente qu'une fraction d'un pour cent de la variance, il change, est partagé entre les actifs
et s'érode sous l'effet de la concurrence.
\end{enumerate}
\end{solution}

\begin{solution}{iq:ml:why-financial-machine-learning-is-different:1}
Pour les actions prises isolément, à fréquence mensuelle, mesuré par rapport à zéro hors échantillon, oui~: les meilleurs modèles publiés atteignent environ
0.4~\%. Vérifier l'indice de référence (zéro ou moyenne), la période, s'il s'agit de la coupe transversale ou du niveau, et la
corrélation qu'elle implique ($\sqrt{0.005} = 0.07$).
\iqlookfor{attentes calibrées du R² en finance, et les questions qui donnent un sens à un nombre.}
\end{solution}

\begin{solution}{iq:ml:why-financial-machine-learning-is-different:2}
L'erreur quadratique espérée en $x$ = bruit + biais au carré + variance. Les rendements sont presque entièrement du bruit, donc la variance coûte cher
et un biais modeste est bon marché~: les modèles régularisés et rigides sont la norme, et les modèles souples exigent beaucoup de données.
\iqlookfor{la décomposition énoncée correctement et appliquée à un problème de faible rapport signal sur bruit.}
\end{solution}

\begin{solution}{iq:ml:why-financial-machine-learning-is-different:3}
Bien moins de dix millions. Pour tout ce qui est commun à toutes les actions (le marché, un facteur), il existe environ 5\,000 jours~; les actions
sont corrélées (au mieux quelques dizaines effectives par jour)~; les étiquettes qui se chevauchent divisent encore le compte~; et, pour un signal lent, les
observations indépendantes sont plus proches du nombre de ses demi-vies dans l'échantillon.
\iqlookfor{dates contre lignes, corrélation en coupe transversale, chevauchement et persistance.}
\end{solution}

\begin{solution}{iq:ml:why-financial-machine-learning-is-different:4}
La moyenne historique du rendement d'une action isolée est une estimation très bruitée~; une prévision qui la bat peut ne l'avoir
fait que parce que l'indice de référence est mauvais~; zéro est une référence fixe et plus exigeante.
\iqlookfor{comprendre que l'indice de référence est un modèle affecté d'une erreur d'estimation.}
\end{solution}

\begin{solution}{iq:ml:why-financial-machine-learning-is-different:5}
\omterm{def:ml:why-financial-machine-learning-is-different:generalisation}{Surapprentissage} (overfitting) dans le backtest (taille de la recherche, fuite d'information), dérive ou décroissance du signal (régime, encombrement (crowding)), écarts d'implémentation
(calendrier des données, univers, variables calculées différemment en production) et le simple bruit (six mois, c'est court). Vérifier
le nombre d'essais du backtest et les statistiques corrigées, le relancer sur les données de la période en production, comparer les variables en production et
hors production, et calculer combien de mois l'écart doit durer pour être significatif.
\iqlookfor{une liste structurée et un test pour chaque élément, y compris la possibilité que rien ne cloche.}
\end{solution}

\begin{solution}{iq:ml:why-financial-machine-learning-is-different:6}
Avec $y$ d'espérance nulle et une prévision $cf$, $R^2(c) = (2c\rho\sigma_y\sigma_f - c^2\sigma_f^2)/\sigma_y^2$, négatif lorsque $c > 2c^\star$
bien que $\rho > 0$. Réduire la prévision (calibrer son échelle sur les données de validation, avec les méthodes du chapitre~11, ou la
calibration isotonique du livre~7), et juger le classement par l'IC.
\iqlookfor{la formule et la différence entre la qualité du classement et la calibration.}
\end{solution}
